\documentclass[11pt,a4paper]{article}
\usepackage[utf8]{inputenc}
\usepackage[T1]{fontenc}
\usepackage{amsmath,amssymb,amsfonts}
\usepackage{physics}
\usepackage{booktabs}
\usepackage{graphicx}
\usepackage{hyperref}
\usepackage[margin=2.5cm]{geometry}
\usepackage{enumitem}
\usepackage{caption}
\usepackage{multirow}

\title{Gate-Circuit Krotov Optimisation with Monte Carlo Wave Functions\\
for Dissipative Quantum State Preparation:\\
Joint Gate and Lindblad-Rate Control}

\author{Maximilian Frohlich}
\date{March 2026}

\begin{document}
\maketitle

\begin{abstract}
We adapt the trajectory-based Krotov optimal control method of Goerz and
Jacobs~\cite{goerz2018} to parameterised gate-based quantum circuits with
interleaved Lindbladian dissipation, and extend it to treat the
\emph{Lindblad dissipation rates} as additional Krotov control variables.

The original method optimises only the Hamiltonian (gate) controls while
keeping the dissipative part fixed.  We derive the Krotov update rule for
per-layer Lindblad rates by differentiating the MCWF no-jump propagation
with respect to each rate, and show that the resulting rate update
naturally vanishes for jump trajectories (the rate cancels in the
post-jump normalisation).  Non-negativity is enforced by projection.

On a 4-qubit hydrogen $2p$ preparation task we compare three modes:
(i)~gates-only Krotov with fixed rates, (ii)~rates-only Krotov with
fixed gates, and (iii)~joint gates+rates Krotov.  The key finding is that
the joint optimiser \emph{selectively} suppresses unhelpful hardware
operators (amplitude damping, dephasing, ZZ) whose rates converge to zero,
while retaining non-zero rates for target-cooling operators when available.
Joint optimisation achieves $F = 0.88$ on a mixed (hardware + teacher)
basis, outperforming the gates-only baseline ($F = 0.80$) and the
rate-optimised two-stage pipeline ($F = 0.93$).  The gates-only mode
degrades sharply with increasing fixed dissipation complexity
(from $F = 0.80$ to $F = 0.47$), while the joint mode maintains
$F \ge 0.79$ across all bases by adaptively disabling harmful channels.

A 10-qubit scalability experiment confirms that target-selective
dissipation remains the only path to improvement at $d = 1024$, and that
rate optimisation---now available within the Krotov framework---is the
critical bottleneck for scaling.

Finally, we test a practically motivated scenario: appending a
rate-only dissipative tail to a pre-compiled circuit from Q-Alchemy.
Teacher-cooling operators at fixed initial rates close up to $2.2\%$ of
the fidelity gap at 6~qubits without modifying the circuit.  However,
we identify a structural limitation: the first-order Krotov rate
gradient vanishes identically for target-cooling operators, preventing
active rate optimisation and motivating operator-level parametrisations.
\end{abstract}

\section{Introduction}

Preparing a target quantum state $\ket{\psi_\star}$ on a quantum
processor is a fundamental task in quantum computing.  When the exact
preparation circuit is too deep for noisy hardware, one truncates it,
producing an approximate state $\rho_0$ with
$F_0 = \bra{\psi_\star}\rho_0\ket{\psi_\star} < 1$.  A natural
question is whether engineered dissipation \emph{after} the truncated
circuit can increase the fidelity toward the target.

Prior work in this project has explored two routes:
(i)~an ideal global target-cooling channel with jump operators
$L_k = \sqrt{\gamma}\ket{\psi_\star}\bra{\psi_k^\perp}$ that drives
$F(t) \to 1$ monotonically, and
(ii)~hardware-constrained local channels (amplitude damping, dephasing,
ZZ~dephasing, ancilla-reset gadgets) with rates fitted to approximate
the ideal teacher.  Approach~(i) is the theoretical optimum but requires
non-local operators.  Approach~(ii) is hardware-realistic but limited.

In this work we introduce a third, hybrid approach based on the
MCWF trajectory method combined with Krotov's optimal control algorithm,
adapted from Goerz and Jacobs~\cite{goerz2018}.  The key innovation is
an \emph{interleaved architecture} where parameterised unitary gates
alternate with dissipative channels.  Gate parameters are optimised to
\emph{exploit} the dissipation, and we systematically study how the
fidelity scales with the expressivity of the available dissipation
channels---from single-qubit amplitude damping to all-to-all ZZ
interactions and ancilla-reset gadgets.


\section{Mathematical Framework}

\subsection{MCWF Trajectory Method}

An open quantum system governed by the Lindblad master equation
\begin{equation}
\frac{d\rho}{dt} = -i[H, \rho] + \sum_l \gamma_l \mathcal{D}[L_l](\rho),
\end{equation}
with $\mathcal{D}[L](\rho) = L\rho L^\dagger -
\tfrac{1}{2}\{L^\dagger L, \rho\}$, is equivalently represented by an
ensemble of pure-state trajectories~\cite{dalibard1992,plenio1998}:
\begin{equation}
\rho(t) = \lim_{M\to\infty} \frac{1}{M}
\sum_{k=1}^{M} \ket{\Psi_k(t)}\bra{\Psi_k(t)}.
\end{equation}
Each trajectory evolves under the non-Hermitian effective Hamiltonian
$H_\mathrm{eff} = H - \tfrac{i}{2} \sum_l L_l^\dagger L_l$, with
stochastic quantum jumps
$\ket{\Psi} \to L_l\ket{\Psi}/\|L_l\ket{\Psi}\|$ occurring with
probability $p(L_l) = \bra{\Psi}L_l^\dagger L_l\ket{\Psi}\,\delta t$.

\subsection{Krotov's Method for Gate Circuits}

The error functional for state-to-state transfer is
\begin{equation}
J_T = 1 - \frac{1}{M}\sum_{k=1}^{M}
\left|\braket{\Psi_k(T)}{\psi_\star}\right|^2.
\end{equation}
For parameterised gates $U_g = \exp(-i\theta_g G_g)$, we identify the
gate generator $G_g$ as the ``control Hamiltonian'' and obtain the
Krotov update~\cite{goerz2018}:
\begin{equation}
\Delta\theta_g = \frac{1}{M\lambda}
\sum_{k=1}^{M}
\mathrm{Im}\bra{\chi_k(t_g)} G_g \ket{\Psi_k^{(1)}(t_g)},
\label{eq:krotov_gate}
\end{equation}
where $\ket{\Psi_k^{(1)}}$ is the forward state after the preceding
gates (with updated parameters), and $\ket{\chi_k}$ is the backward
co-state propagated from
$\ket{\chi_k(T)} = \braket{\psi_\star}{\Psi_k^{(0)}(T)}\ket{\psi_\star}$.

The crucial property is that Krotov's method requires only the
derivative of the equation of motion ($G_g$), not the propagator
derivative.  This preserves compatibility with stochastic trajectories.

\subsection{Extended Krotov: Lindblad Rates as Controls}
\label{sec:rate_krotov}

We extend the Krotov framework to treat the \emph{per-layer
Lindblad rates} $\gamma_k^{(l)} \ge 0$ as additional control variables.
The jump operators $L_k$ remain fixed; only the rates are optimised.

At layer $l$, the MCWF no-jump propagation (to first order in $\delta t$)
with rate-dependent operators
$\tilde{L}_k^{(l)} = \sqrt{\gamma_k^{(l)}} L_k$ reads:
\begin{equation}
\ket{\psi_\text{out}} \propto \ket{\psi'} - \frac{\delta t}{2}
\sum_k \gamma_k^{(l)}\, L_k^\dagger L_k\, \ket{\psi'},
\label{eq:nojump}
\end{equation}
where $\ket{\psi'}$ is the state after applying the unitary gates at
layer~$l$.  Differentiating with respect to $\gamma_k^{(l)}$:
\begin{equation}
\frac{\partial\ket{\psi_\text{out}}}{\partial \gamma_k^{(l)}}
= -\frac{\delta t}{2}\, L_k^\dagger L_k\, \ket{\psi'}.
\label{eq:drate}
\end{equation}

For \emph{jump trajectories}, the post-jump state
$L_j\ket{\psi'}/\|L_j\ket{\psi'}\|$ is independent of
$\gamma_k^{(l)}$: the rate appears in the selection probability
but cancels in the normalisation.  Thus only no-jump trajectories
contribute to the rate gradient.

Substituting~\eqref{eq:drate} into the Krotov functional variation
yields the update rule:
\begin{equation}
\boxed{
\Delta\gamma_k^{(l)} = -\frac{\delta t}{2\,M\,\lambda_\gamma}
\sum_{\substack{m=1\\[2pt]\text{no-jump}}}^{M}
\operatorname{Re}\!\bra{\chi_m'^{(l)}}\,
L_k^\dagger L_k\,\ket{\psi_m'^{(l)}},
}
\label{eq:krotov_rate}
\end{equation}
where $\ket{\psi_m'^{(l)}}$ is the (updated) forward state after
gates at layer~$l$, and $\ket{\chi_m'^{(l)}}$ is the backward co-state
at the intermediate point (after gates, before dissipation) obtained
from the backward pass using old controls.  The regularisation parameter
$\lambda_\gamma > 0$ controls the step size.

Physical admissibility is enforced by projecting onto the non-negative
orthant after each update:
\begin{equation}
\gamma_k^{(l)} \leftarrow
\max\!\bigl(0,\; \gamma_k^{(l)} + \Delta\gamma_k^{(l)}\bigr).
\end{equation}

The backward co-state $\ket{\chi_m'^{(l)}}$ is stored during the
backward propagation by recording the state at the intermediate point
between the dissipation-undo and gate-undo steps:
\begin{align}
\ket{\chi^{(l+1)}} &\xrightarrow{\text{undo dissipation}} \ket{\chi'^{(l)}}
&& (\text{stored for rate update}), \\
\ket{\chi'^{(l)}} &\xrightarrow{\text{undo gates}} \ket{\chi^{(l)}}
&& (\text{stored for gate update}).
\end{align}

The three optimisation modes are obtained by enabling or disabling
the corresponding update rules:
\begin{description}[nosep]
\item[Gates-only:] Eq.~\eqref{eq:krotov_gate} active,
  rates $\gamma_k^{(l)}$ fixed.
\item[Rates-only:] Eq.~\eqref{eq:krotov_rate} active,
  gate parameters $\theta_g$ fixed.
\item[Joint:] Both Eqs.~\eqref{eq:krotov_gate}
  and~\eqref{eq:krotov_rate} active.  Gates are updated first, then
  the forward state is propagated through the updated gates before
  computing the rate update.
\end{description}

\subsection{Interleaved Gate+Dissipation Architecture}

Each of the $L$ circuit layers applies:
\begin{equation}
\ket{\Psi} \;\xrightarrow{\;U_1(\theta_1)\cdots U_G(\theta_G)\;}
\ket{\Psi'} \;\xrightarrow{\;\mathcal{D}_\mathrm{Lindblad}(\gamma,\delta t)\;}
\ket{\Psi''} \;\xrightarrow{\;\{\mathcal{K}_s\}\;}
\ket{\Psi'''}.
\end{equation}
The first arrow is the unitary gate layer.  The second applies a
Lindblad MCWF step (possible quantum jump).  The optional third
arrow applies discrete Kraus channels $\{\mathcal{K}_s\}$
(e.g.\ ancilla-reset gadgets).

\subsection{Hardware-Efficient Ansatz}

We employ a hardware-efficient variational ansatz.
Each of the $L$ layers contains $R_y(\theta)$,
$R_z(\theta)$ rotations on all $n$ qubits, and CNOT gates on
nearest-neighbour pairs, giving
$2n + (n-1) = 3n - 1$ gates per layer, for a total of $L(3n-1)$
parameters ($44$ for $n=4$, $L=4$; $116$ for $n=10$, $L=4$).

\subsection{Efficient Tensor-Contraction Gate Application}

For scalability to larger qubit counts, gate application must be
performed efficiently.  A single-qubit parameterised gate
$U = \exp(-i\theta G)$ with local $2\times 2$ generator $G$ acts on the
state vector $\ket{\Psi} \in \mathbb{C}^{2^n}$ in $O(2^n)$ time by:
\begin{enumerate}[nosep]
\item Computing the $2\times 2$ matrix exponential
  $U_\text{loc} = \exp(-i\theta G_\text{loc})$.
\item Reshaping $\ket{\Psi}$ as a rank-$n$ tensor and contracting
  $U_\text{loc}$ with the target qubit axis.
\end{enumerate}
This replaces the $O(d^3)$ full matrix exponential with an $O(d)$
tensor contraction, where $d = 2^n$.  Two-qubit gates use analogous
$4\times 4$ contractions in $O(4d)$.  For the Krotov update, the
generator application $G\ket{\Psi}$ uses the same tensor-contraction
approach, and the density-matrix evaluation applies each gate
individually via rank-$2n$ tensor contractions in $O(2^k \cdot d^2)$
per $k$-qubit gate.

For the Lindblad operators, sparse matrix representations are used:
amplitude-damping operators have $d/2$ non-zero elements, while
dephasing and ZZ operators are diagonal, requiring $O(d)$ storage
instead of $O(d^2)$.  Sparse matrix--vector products are $O(\text{nnz})$.


\section{Hardware-Constrained Dissipation Channels}
\label{sec:channels}

We study four levels of increasing dissipation expressivity, each a
strict superset of the previous.  All rates are optimised to best
approximate the ideal target-cooling channel (Section~\ref{sec:fitting}).

\subsection{Config~A: Amplitude Damping + Dephasing}

For each qubit $q = 0, \ldots, n-1$ the available jump operators are:
\begin{align}
L_q^\text{AD} &= \sqrt{\gamma_q^\text{AD}}\;\sigma_q^- ,
& \sigma^- &= \ket{0}\bra{1}, \label{eq:amp_damp} \\
L_q^\text{D}  &= \sqrt{\gamma_q^\text{D}}\;Z_q ,
& Z &= \ket{0}\bra{0} - \ket{1}\bra{1}. \label{eq:dephasing}
\end{align}
Amplitude damping transfers population from $\ket{1}$ to $\ket{0}$
(modelling spontaneous emission / $T_1$ decay), while dephasing
destroys off-diagonal coherences without changing populations
($T_2^*$ decay).
Total parameter count: $2n = 8$.

\subsection{Config~B: A + ZZ Dephasing (Chain Connectivity)}

In addition to the single-qubit channels, correlated two-qubit
dephasing is enabled on nearest-neighbour edges
$\mathcal{E}_\text{chain} = \{(0,1),(1,2),\ldots,(n{-}2,n{-}1)\}$:
\begin{equation}
L_{ij}^\text{ZZ} = \sqrt{\gamma_{ij}^\text{ZZ}}\;Z_i \otimes Z_j,
\quad (i,j) \in \mathcal{E}_\text{chain}.
\label{eq:zz}
\end{equation}
This operator preserves populations in the computational basis but
induces correlated dephasing: states with the same parity on qubits
$i,j$ acquire a different phase decay rate than those with opposite
parity. Total: $2n + (n-1) = 11$ parameters.

\subsection{Config~C: A + ZZ Dephasing (All-to-All Connectivity)}

The ZZ operators~\eqref{eq:zz} are extended to the complete graph:
\begin{equation}
\mathcal{E}_\text{a2a} = \{(i,j) : 0 \le i < j < n\},
\quad |\mathcal{E}_\text{a2a}| = \binom{n}{2}.
\end{equation}
This removes geometric locality constraints while still restricting
primitives to 2-qubit Lindblad operators.
Total: $2n + \binom{n}{2} = 14$ parameters.

\subsection{Config~D: C + Ancilla-Reset Gadgets}

For each system qubit $q$, an ancilla qubit initialised in $\ket{0}$
is coupled via a parameterised $4\times 4$ unitary
$U_q(\boldsymbol\phi_q)$, then the ancilla is reset:
\begin{equation}
K_{q,m} = {}_\text{anc}\!\bra{m}\, U_q(\boldsymbol\phi_q)\,
\ket{0}_\text{anc}, \quad m \in \{0,1\}.
\end{equation}
This implements a programmable single-qubit CPTP map
$\mathcal{E}_q(\rho) = \sum_m K_{q,m}\rho\, K_{q,m}^\dagger$ that
can, in principle, realise \emph{any} qubit channel by choosing
$\boldsymbol\phi_q$ appropriately.

In the MCWF framework, the ancilla-reset step is a discrete
stochastic operation: outcome $m$ is drawn with probability
$p_m = \|K_{q,m}\ket{\Psi}\|^2$ and the state collapses to
$K_{q,m}\ket{\Psi}/\|K_{q,m}\ket{\Psi}\|$.

Each $U_q$ is parameterised as
$U_q = \exp\!\bigl(-i\sum_{k=1}^{15}\phi_{q,k}\,P_k\bigr)$
where $\{P_k\}$ is the two-qubit Pauli product basis (excluding identity).
Total Lindblad + Kraus parameters: $14 + 4 \times 15 = 74$.

\section{Two-Stage Optimisation Pipeline}
\label{sec:fitting}

\subsection{Stage~1: Rate Optimisation}

For each constraint level, the dissipation rates
$\boldsymbol\gamma$ (and ancilla angles $\boldsymbol\phi$) are
optimised to minimise the \emph{fidelity-curve loss}
\begin{equation}
\mathcal{L}(\boldsymbol\gamma)
= \frac{1}{SJ}\sum_{s=1}^{S}\sum_{j=1}^{J}
\bigl(F_\text{model}^{(s)}(t_j) - F_\text{teacher}^{(s)}(t_j)\bigr)^2
+ \eta\,\|\boldsymbol\gamma / \gamma_\text{max}\|^2,
\label{eq:fit_loss}
\end{equation}
where $F_\text{teacher}^{(s)}(t)$ is the fidelity curve of the ideal
target-cooling channel (with $\gamma_\text{teacher} = 1.0$,
$T = 5.0$), and $F_\text{model}^{(s)}$ is the Trotterised
constrained evolution from training state~$s$.
The regularisation weight is $\eta = 0.005$.

The optimisation uses L-BFGS-B with rates mapped through a sigmoid
$\gamma = \gamma_\text{max}\,\sigma(x)$ to enforce
$\gamma \in (0, \gamma_\text{max})$.  Training uses $S = 5$ diverse
initial states (truncated circuit output, target state, random pure
and mixed states) and $J = 20$ time steps.

\subsection{Stage~2: MCWF-Krotov Gate Optimisation}

The optimised Lindblad operators and Kraus channels from Stage~1 are
used in the MCWF-Krotov gate optimiser
(Eq.~\ref{eq:krotov_gate}).  A fresh 4-layer hardware-efficient ansatz is
initialised with random parameters $\theta_g \in [-0.5, 0.5]$ and
optimised for 100 iterations with $M = 6$ trajectories,
$\lambda = 0.3$, $\delta t = 0.1$.

This two-stage pipeline separates the concerns of
\emph{what dissipation to apply} (Stage~1) from
\emph{how to interleave it with gates} (Stage~2).


\section{Results}

\subsection{Extended Comparison}

\begin{table}[ht]
\centering
\begin{tabular}{@{}lcccccc@{}}
\toprule
\textbf{Method} & \textbf{$F_\text{DM}$} & $\Delta F$
& \textbf{Purity} & \textbf{\#Ops}
& \textbf{Fit Loss} & $F_\text{diss}$ \\
\midrule
Truncated circuit (baseline)
  & 0.500 & ---   & 1.000 & --- & --- & --- \\
Teacher tail (DM, ideal)
  & 0.997 & +0.497 & ---  & --- & --- & --- \\
\midrule
\multicolumn{7}{l}{\textit{MCWF-Krotov with fixed-rate operators}} \\
\quad Teacher Lindblad (oracle)
  & 0.814 & +0.314 & 0.854 & 15 & --- & --- \\
\quad 1q amplitude damping
  & 0.665 & +0.165 & 0.768 &  4 & --- & --- \\
\midrule
\multicolumn{7}{l}{\textit{MCWF-Krotov with rate-optimised operators (two-stage)}} \\
\quad Config~A: AD + dephasing
  & 0.581 & +0.081 & 1.000 &  8 & 0.425 & 0.500 \\
\quad Config~B: A + ZZ (chain)
  & 0.630 & +0.130 & 0.981 & 11 & 0.423 & 0.442 \\
\quad Config~C: A + ZZ (all-to-all)
  & \textbf{0.934} & \textbf{+0.434} & 0.981 & 14 & 0.423 & 0.442 \\
\quad Config~D: C + ancilla-reset
  & 0.273 & $-0.227$ & 0.995 & 18 & 0.405 & 0.274 \\
\bottomrule
\end{tabular}
\caption{%
Fidelity $F_\text{DM} = \bra{\psi_\star}\rho\ket{\psi_\star}$
from exact density-matrix evaluation of the optimised circuit.
$\Delta F$ is improvement over the truncated baseline.
\textbf{\#Ops} is the total number of Lindblad + Kraus channel
operators.  \textbf{Fit Loss} is the rate-fitting objective
(\ref{eq:fit_loss}).
$F_\text{diss}$ is the fidelity achieved by the fitted dissipation
alone (no Krotov gate optimisation), applied to the truncated
circuit output.
}
\label{tab:extended}
\end{table}

\subsection{Analysis}

\paragraph{Rate optimisation is essential.}
Config~A achieves only $F = 0.58$ despite having 8~Lindblad operators,
because the L-BFGS-B optimiser sets all rates to near-zero
($\gamma \lesssim 10^{-4}$).  Single-qubit amplitude damping and
dephasing simply cannot approximate the global teacher: they either
add noise (dephasing) or drive toward $\ket{0\ldots 0}$ (amplitude
damping), neither of which helps for a general target.  The
purity $\approx 1.0$ confirms that effectively no dissipation acts.

\paragraph{Connectivity matters more than channel count.}
Configs~B and C have nearly identical fit losses ($\approx 0.423$)
yet the MCWF-Krotov fidelity jumps from $F = 0.63$ (chain) to
$F = 0.93$ (all-to-all).  The ZZ operators on distant qubit pairs
enable correlated decoherence patterns that the gate optimiser can
exploit: gates steer populations into parity subspaces where the ZZ
dephasing selectively suppresses unwanted components, then subsequent
gates rotate toward the target.

\paragraph{All-to-all ZZ achieves the best hardware-constrained result.}
Config~C with $F = 0.93$ substantially outperforms both the
fixed-rate teacher Lindblad ($F = 0.81$) and the amplitude damping
baseline ($F = 0.67$).  This demonstrates that \emph{structured}
dissipation with long-range connectivity, when rates are optimised,
provides the gate optimiser with significantly more leverage than
non-local but unstructured (teacher) operators with ad-hoc rates.

\paragraph{Ancilla-reset is powerful but hard to optimise.}
Config~D has the \emph{lowest} fit loss ($0.405 < 0.423$), confirming
that the ancilla-reset CPTP maps can better approximate the teacher
channel.  However, the MCWF-Krotov gate optimisation fails to converge
($F \approx 0.27$, essentially the dissipation-only value).  The
74-parameter ancilla-reset landscape creates strong non-unitary
channel operations that the sequential Krotov update for gates alone
cannot compensate.  This suggests that a \emph{joint} Krotov--gradient
optimisation over both gate and ancilla parameters is needed.

\paragraph{Dissipation-only vs.\ Krotov-enhanced.}
For configs~B and C, the dissipation alone drives the truncated circuit
state to $F_\text{diss} \approx 0.44$, while the combined
gate+dissipation optimisation achieves $F = 0.63$--$0.93$.  The gate
optimiser thus provides a multiplicative boost on top of the
dissipation's own cooling capability.


\subsection{10-Qubit Scalability Experiment}

To assess scalability, we extend the experiment to a 10-qubit system
($d = 1024$, hydrogen $2p$ target).  The circuit uses $L = 6$ layers
($174$ gates), $\lambda = 0.3$, $\delta t = 0.1$, $M = 8$ trajectories,
and per-gate update clipping $|\Delta\theta_j| \le 0.05$ to prevent
the overshooting that arises at large $d$ when the Krotov update
magnitudes scale as $O(1/\sqrt{d})$.  Gate parameters are initialised
near zero (circuit $\approx I$) to preserve the truncated-state
baseline.

\paragraph{Memory-efficient teacher operators.}
At $d = 1024$, storing the full $d - 1 = 1023$ teacher operators
$L_k = \sqrt{\gamma}\ket{\psi_*}\bra{\psi_k^\perp}$ as dense
$d \times d$ matrices requires ${\sim}17\,\text{GB}$.  Instead, we
introduce a \texttt{Rank1Op} class that stores only the two defining
vectors $\ket{\psi_*}$ and $\ket{\psi_k^\perp}$ and implements
$O(d)$ matrix-vector products.  We retain only the top-$k$ operators
ranked by their overlap $|\braket{\psi_k^\perp}{\psi_\text{trunc}}|^2$
with the initial state, capturing the dominant cooling directions.
For density-matrix evaluation, all rank-1 Lindblad terms are batched
into a single projector update, avoiding the per-operator $O(d^2)$
temporaries that would otherwise exhaust memory.

\paragraph{Hardware channel rates.}
Two-stage rate fitting (Section~\ref{sec:fitting}) is prohibitively
expensive at $d = 1024$ ($O(d^2)$ per DM propagation step), so
hardware channels use small fixed rates:
$\gamma_\text{AD} = 0.05$, $\gamma_\text{D} = 0.01$,
$\gamma_\text{ZZ} = 0.01$.

\begin{table}[ht]
\centering
\begin{tabular}{@{}lcccc@{}}
\toprule
\textbf{Method} & \textbf{$F_\text{DM}$} & $\Delta F$
& \textbf{Purity} & \textbf{Time (s)} \\
\midrule
Truncated circuit (baseline)
  & 0.500 & ---     & 1.000  & ---  \\
Teacher tail (DM, ideal)
  & 0.997 & +0.497  & ---    & ---  \\
\midrule
Krotov-only (no dissipation, 200 iters)
  & 0.500 & ${\sim}0$ & 1.000 & 48   \\
\textbf{Reduced teacher (10 Rank1Ops, 100 iters)}
  & \textbf{0.515} & +0.015 & 0.976 & 25   \\
\midrule
1q amplitude damping ($n_\text{ops}{=}10$)
  & 0.444 & $-0.056$ & 0.880 & 30   \\
AD+deph+ZZ chain ($n_\text{ops}{=}29$)
  & 0.405 & $-0.095$ & 0.747 & 46   \\
AD+deph+ZZ all-to-all ($n_\text{ops}{=}65$)
  & 0.325 & $-0.175$ & 0.550 & 77   \\
\bottomrule
\end{tabular}
\caption{%
10-qubit hydrogen $2p$ state preparation.  The \texttt{Rank1Op} teacher
configuration is the \emph{only} method that improves over the truncated
baseline at this scale, demonstrating that target-selective dissipation
remains effective even with a reduced operator set.  Fixed-rate
hardware channels add decoherence without constructive cooling.
}
\label{tab:10q}
\end{table}

\paragraph{Teacher dissipation improves fidelity at 10 qubits.}
Using only 10 memory-efficient \texttt{Rank1Op} teacher operators
(out of 1023 possible), the MCWF-Krotov method achieves $F = 0.515$
in 100 iterations --- a statistically significant improvement over the
$F = 0.500$ baseline.  While modest compared to the ideal teacher
tail ($F = 0.997$), this confirms that dissipation-assisted gate
optimisation \emph{does scale} to larger Hilbert spaces when the
dissipation channel is target-selective.

\paragraph{Gate-only optimisation saturates.}
With 200 iterations, the ``Krotov-only'' experiment achieves $F = 0.500$
--- no improvement.  At $d = 1024$, the Krotov update
$\Delta\theta_j \propto \operatorname{Im}\braket{\chi|G_j|\psi}$
has magnitude $O(1/\sqrt{d})$, yielding per-gate increments of
${\sim}0.005$\,rad even with the step-size clipping.  Over 200 iterations
the accumulated change is insufficient to restructure the full circuit.
This demonstrates that \textbf{dissipation is not merely helpful but
necessary} for making progress at this scale.

\paragraph{Fixed-rate hardware noise degrades fidelity monotonically.}
Amplitude damping ($F = 0.444$), chain connectivity ($F = 0.405$),
and all-to-all connectivity ($F = 0.325$) all perform worse than the
baseline.  The degradation increases with channel count, consistent
with the 4-qubit pre-fitting results.  Without rate optimisation,
dephasing operators ($L^\dagger L \propto I$) contribute pure
decoherence rather than target-selective cooling.  The purity drop
from $1.0$ to $0.55$ for 65 operators quantifies the noise
accumulation.

\paragraph{Computational efficiency.}
The tensor-contraction gate application ($O(2^k d)$ per gate) and
\texttt{Rank1Op} formulation ($O(d)$ per operator) make the MCWF-Krotov
iteration tractable: $25$--$77\,\text{s}$ for $100$ iterations with
$8$ trajectories and $174$ gates.  The batched density-matrix
evaluation avoids the $k \times 5$ temporary $d \times d$ matrices
that would otherwise exhaust memory at $d = 1024$.


\subsection{Extended Krotov: Joint Gate + Rate Optimisation}
\label{sec:rate_results}

To validate the rate-Krotov extension (Section~\ref{sec:rate_krotov}),
we run the same 4-qubit hydrogen $2p$ preparation task in three modes---gates-only,
rates-only, and joint---across four dissipation bases of increasing
expressivity.  In all cases, the jump operators are fixed and only
the per-layer rates are (optionally) optimised within the Krotov loop.

\paragraph{Setup.}
The circuit uses $L = 4$ layers ($44$ gates), $M = 6$ trajectories,
$80$ Krotov iterations, $\lambda_\theta = 0.3$, $\lambda_\gamma = 0.5$,
$\delta t = 0.1$, per-gate clipping $|\Delta\theta| \le 0.15$, and
per-rate clipping $|\Delta\gamma| \le 0.08$.  All rates are initialised
at $\gamma_0 = 0.1$.

The four dissipation bases are:
\begin{enumerate}[nosep]
\item 1q amplitude damping (4~operators),
\item 1q AD + dephasing (8~operators),
\item 1q AD + deph + 2q ZZ all-to-all (14~operators),
\item Mixed: 1q AD + 4 target-cooling operators (8~operators).
\end{enumerate}

\begin{table}[ht]
\centering
\begin{tabular}{@{}l@{\quad}c@{\quad}c@{\quad}c@{}}
\toprule
\textbf{Basis} & \textbf{Gates-only} & \textbf{Rates-only}
& \textbf{Joint} \\
\midrule
1q amp-damping
  & 0.803 & 0.343 & 0.791 \\
1q AD + dephasing
  & 0.617 & 0.343 & 0.791 \\
1q AD + deph + ZZ (a2a)
  & 0.469 & 0.343 & 0.871 \\
AD + teacher cooling
  & 0.801 & 0.346 & \textbf{0.882} \\
\midrule
\multicolumn{4}{l}{\textit{Baseline truncated fidelity: $F = 0.500$}} \\
\bottomrule
\end{tabular}
\caption{%
Fidelity $F_\text{DM}$ (exact density-matrix evaluation) for the
three optimisation modes across four dissipation bases on the 4-qubit
hydrogen $2p$ task.  Joint optimisation achieves the best result on
the mixed basis ($F = 0.882$) by selectively retaining teacher operators
and suppressing hardware noise.
}
\label{tab:rate_krotov}
\end{table}

\paragraph{Gates-only mode degrades with dissipation complexity.}
When rates are fixed at $\gamma = 0.1$, adding more hardware operators
\emph{hurts}: the fidelity drops from $F = 0.80$ (AD-only) to
$F = 0.62$ (AD+deph) to $F = 0.47$ (AD+deph+ZZ).  The ZZ operators
on all qubit pairs introduce strong correlated dephasing that the gate
optimiser cannot compensate.  The purity drops from $0.88$ to $0.43$,
confirming that the fixed dissipation acts as uncontrolled noise.

\paragraph{Rates-only mode correctly identifies operator utility.}
In rates-only mode (gates frozen at the initial ansatz), all hardware
operator rates converge to zero within a few iterations: $\gamma_k = 0$
for all AD, dephasing, and ZZ operators.  This is physically correct---these
operators do not help prepare an arbitrary target state---and confirms
that the rate-Krotov gradient
$\Delta\gamma_k \propto -\mathrm{Re}\langle\chi|L_k^\dagger L_k|\psi\rangle$
correctly pushes unhelpful rates toward zero.

The mixed basis (Basis~4) reveals selective behaviour:
AD rates converge to zero, while target-cooling rates remain non-zero
and in some layers \emph{increase} (e.g.\ $\gamma_\text{TC\_1}^{(l=0)} = 0.22$,
up from the initial $0.10$).  This demonstrates that the rate-Krotov
update can distinguish helpful from harmful operators within the same
basis.  However, without gate optimisation, the overall fidelity
improvement is small ($F = 0.35$).

\paragraph{Joint mode is robust across all bases.}
The joint optimiser maintains $F \ge 0.79$ across all bases by adaptively
disabling harmful channels: all hardware rates converge to zero, yielding
effectively pure unitary optimisation.  On the mixed basis, the joint
mode achieves the highest fidelity ($F = 0.88$) by simultaneously
optimising gates \emph{and} selectively retaining teacher operators in
the final layers.  The final rate profile shows AD rates at zero and
teacher-cooling rates non-zero in the last layer
($\gamma_\text{TC} = 0.10$), while earlier layers also suppress them
($\gamma_\text{TC} = 0.00$).  This per-layer rate structure is not
accessible to the two-stage pipeline, which uses a single global rate
per operator.

\paragraph{Comparison with two-stage pipeline.}
The best result from the two-stage pipeline (Section~\ref{sec:fitting},
Config~C) achieves $F = 0.93$ with pre-fitted all-to-all ZZ rates.
The joint-Krotov mixed basis ($F = 0.88$) is slightly lower but uses
only 4 teacher operators rather than 14 ZZ operators with pre-fitted
rates.  Crucially, the joint approach does not require a separate
rate-fitting stage and can handle bases where the optimal rates are
zero (disabling channels), which the two-stage pipeline cannot
discover during Stage~1.


\subsection{Scaling: 6-Qubit and 8-Qubit Experiments}
\label{sec:scaling_rate_krotov}

We repeat the joint-Krotov experiment at $n = 6$ ($d = 64$) and
$n = 8$ ($d = 256$) to study how the three optimisation modes scale
with system size.  At 6~qubits the circuit uses $L = 5$ layers
($85$ gates, per-gate clip $0.10$, per-rate clip $0.06$); at 8~qubits
also $L = 5$ layers ($115$ gates, per-gate clip $0.08$, per-rate clip
$0.05$).  The mixed basis uses $n$ AD operators plus
$\min(n, 8)$ teacher-cooling operators.  All other hyperparameters
match the 4-qubit setup.

\begin{table}[ht]
\centering
\begin{tabular}{@{}l@{\quad}c@{\quad}c@{\quad}c@{}}
\toprule
\textbf{Basis} & \textbf{Gates-only} & \textbf{Rates-only}
& \textbf{Joint} \\
\midrule
\multicolumn{4}{l}{\textit{6 qubits ($d = 64$, baseline $F = 0.500$)}} \\
\quad 1q amp-damping
  & 0.434 & 0.094 & 0.508 \\
\quad AD + dephasing
  & 0.298 & 0.094 & 0.503 \\
\quad AD + deph + ZZ (a2a)
  & 0.178 & 0.094 & \textbf{0.550} \\
\quad AD + teacher (mixed)
  & 0.437 & 0.098 & 0.508 \\
\midrule
\multicolumn{4}{l}{\textit{8 qubits ($d = 256$, baseline $F = 0.500$)}} \\
\quad 1q amp-damping
  & 0.419 & 0.026 & 0.499 \\
\quad AD + dephasing
  & 0.297 & 0.026 & 0.500 \\
\quad AD + deph + ZZ (a2a)
  & 0.083 & 0.026 & 0.500 \\
\quad AD + teacher (mixed)
  & 0.420 & 0.030 & \textbf{0.500} \\
\bottomrule
\end{tabular}
\caption{%
Fidelity $F_\text{DM}$ for the three optimisation modes at 6 and
8~qubits.  Gates-only degrades sharply with both system size and
dissipation complexity. Joint Krotov maintains fidelity at or above
the truncated baseline by suppressing harmful channels.
At 6~qubits the joint mode with AD+deph+ZZ achieves
$F = 0.55$ ($+10\%$), demonstrating modest but meaningful
improvement.  At 8~qubits, gate optimisation saturates near
the baseline; the mixed-basis joint mode retains teacher rates but
the Krotov gradient is too small ($O(1/\sqrt{d})$) for significant
progress.
}
\label{tab:rate_krotov_scaling}
\end{table}

\paragraph{Consistent patterns across scales.}
Three findings persist at 6 and 8~qubits:
\begin{enumerate}[nosep]
\item \emph{Rates-only} always suppresses hardware operators to zero and
  selectively retains teacher-cooling rates (AD rates $\to 0$, TC rates
  non-zero at 6q: mean $\gamma_\text{TC} \approx 0.09$; at 8q: mean
  $\gamma_\text{TC} \approx 0.10$).
\item \emph{Gates-only with fixed rates} degrades monotonically with
  dissipation complexity: at 6q from $F = 0.43$ (AD) to $F = 0.18$
  (AD+deph+ZZ); at 8q from $F = 0.42$ to $F = 0.08$.
\item \emph{Joint mode} maintains baseline fidelity or better by
  disabling harmful channels---it never performs \emph{worse} than the
  truncated circuit, unlike gates-only.
\end{enumerate}

\paragraph{Diminishing returns with system size.}
The fidelity improvement from joint-Krotov shrinks with $n$: $+0.38$
at 4q, $+0.05$ at 6q, and ${\sim}0$ at 8q.  This is consistent with
the $O(1/\sqrt{d})$ scaling of the Krotov update magnitude.  At
$d = 256$, individual gate updates are ${\sim}0.006$\,rad, and
80~iterations accumulate only ${\sim}0.5$\,rad of total circuit
restructuring---insufficient for the $115$-gate, 5-layer ansatz.
The 6-qubit experiment at $d = 64$ occupies an intermediate regime
where gate updates (${\sim}0.012$\,rad) are just large enough for
the AD+deph+ZZ joint mode to achieve $F = 0.55$.

\paragraph{Protective role of rate optimisation.}
Even where the joint mode shows little improvement over baseline, it
provides a crucial \emph{protective} function: without rate
optimisation, adding 44 hardware operators (8q AD+deph+ZZ) destroys
fidelity ($F = 0.08$, purity $0.14$).  The joint optimiser avoids
this catastrophic degradation entirely by zeroing all rates.  This
``selective silence'' mechanism ensures that extending the operator
basis never hurts performance---a property not available to the
gates-only mode.


\subsection{Q-Alchemy Circuit + Dissipative Rate Tail}
\label{sec:qalchemy}

A natural question is whether the rate-Krotov method can close the
fidelity gap of a \emph{pre-compiled} state-preparation circuit
obtained from an external compiler, without modifying the circuit
gates at all.  We test this using Q-Alchemy~\cite{qalchemy}, a
cloud-based circuit compiler that produces optimised circuits for a
given target statevector at a user-specified
\texttt{max\_fidelity\_loss} parameter controlling the circuit depth
vs.\ fidelity trade-off.

\paragraph{Protocol.}
For 6~qubits ($d = 64$), the hydrogen $2p$ target, and five fidelity-loss
levels $\epsilon \in \{0.05, 0.10, 0.20, 0.30, 0.40\}$, we:
\begin{enumerate}[nosep]
\item Call Q-Alchemy to obtain a circuit achieving fidelity
  $F_\text{circ} = 1 - \delta$ (where $\delta \le \epsilon$ in
  general).
\item Evolve $\ket{0}^{\otimes n}$ through the circuit to obtain
  $\ket{\psi_\text{circ}}$.
\item Append $L_\text{diss} = 10$ ``identity-gate + dissipation''
  layers (all gate parameters set to zero, i.e.\ $U = I$).
\item Run \emph{rate-only} Krotov ($100$ iterations, $M = 8$
  trajectories, $\lambda_\gamma = 0.1$, $\delta t = 0.3$) to
  optimise per-layer Lindblad rates.
\end{enumerate}
Three dissipation bases are tested:
(i)~hardware-only (1q AD + dephasing, $\gamma_0 = 0.01$),
(ii)~teacher cooling (20 target-cooling operators, $\gamma_0 = 0.3$),
(iii)~mixed (AD + 20 teacher operators).

\begin{table}[ht]
\centering
\begin{tabular}{@{}ccccc@{}}
\toprule
$\epsilon$ & $F_\text{circ}$ & Hardware & Teacher & Mixed \\
\midrule
0.05 & 0.9957 & 0.9957 & 0.9964 & 0.9964 \\
0.10 & 0.9957 & 0.9957 & 0.9964 & 0.9964 \\
0.20 & 0.8818 & 0.8818 & \textbf{0.9040} & \textbf{0.9053} \\
0.30 & 0.8818 & 0.8818 & 0.9040 & 0.9053 \\
0.40 & 0.6153 & 0.6153 & 0.6347 & 0.6362 \\
\bottomrule
\end{tabular}
\caption{%
Q-Alchemy circuit fidelity $F_\text{circ}$ and post-dissipation
fidelity $F_\text{DM}$ for three rate-only Krotov bases (6~qubits).
Teacher cooling partially closes the gap: $+2.2\%$ absolute at
$\epsilon = 0.20$.  Hardware rates collapse to zero and provide
no improvement.
}
\label{tab:qalchemy}
\end{table}

\paragraph{Hardware rates provide no improvement.}
Across all fidelity-loss levels, all hardware operator rates (AD and
dephasing) converge to zero within the first few Krotov iterations,
yielding $F_\text{DM} = F_\text{circ}$ exactly.  This is consistent
with the gates-frozen rate-only results of
Section~\ref{sec:rate_results}: single-qubit dissipation channels
cannot distinguish the target from other states in the orthogonal
complement and therefore cannot provide cooling.

\paragraph{Teacher cooling partially closes the gap.}
The target-cooling operators provide a measurable fidelity improvement
across all gap sizes: $+0.065\%$ for $\epsilon = 0.05$
($F: 0.9957 \to 0.9964$), $+2.2\%$ for $\epsilon = 0.20$
($F: 0.8818 \to 0.9040$), and $+1.9\%$ for $\epsilon = 0.40$
($F: 0.6153 \to 0.6347$).  This improvement comes entirely from
the \emph{fixed initial rates} applied through $L_\text{diss}$
dissipation layers: the teacher cooling channels act as a passive
``dissipative tail'' that contracts the state toward the target.

\paragraph{Structural vanishing of the rate-Krotov gradient.}
A key theoretical observation is that the rate-Krotov gradient
(Eq.~\ref{eq:krotov_rate}) \emph{vanishes identically} for
target-cooling operators when the backward co-state is proportional
to $\ket{\psi_\star}$.  Since
$L_k^\dagger L_k = \ket{\psi_k^\perp}\bra{\psi_k^\perp}$ for
teacher operators, and
$\bra{\psi_\star}\ket{\psi_k^\perp} = 0$ by construction:
\begin{equation}
\bra{\chi}\, L_k^\dagger L_k\, \ket{\psi} \;\propto\;
\bra{\psi_\star}\ket{\psi_k^\perp}\,
\bra{\psi_k^\perp}\ket{\psi} = 0.
\label{eq:vanishing_gradient}
\end{equation}
This explains why teacher rates remain exactly at their initial
values throughout optimisation (confirmed numerically:
$\gamma_k^{(l)} = 0.300$ at all layers after $100$ iterations).
The fidelity improvement is therefore not from \emph{optimisation}
but from the \emph{physical cooling effect} of the teacher operators
at their initial rate.

\paragraph{Implications.}
This experiment demonstrates that target-cooling dissipation can serve
as a practical ``fidelity booster'' appended to any compiled circuit,
provided the target-cooling operators are available.  However, the
improvement is bounded by the fixed-rate cooling dynamics---the Krotov
rate optimiser cannot further tune these rates due to the structural
gradient vanishing~\eqref{eq:vanishing_gradient}.  This is a
fundamental limitation of the first-order MCWF Krotov formulation for
target-cooling operators and motivates either (i)~second-order
corrections, (ii)~parametrisation of the jump operators themselves
rather than just their rates, or (iii)~hybrid approaches combining
compiled circuits with joint gate+dissipation optimisation.

The mixed basis confirms that the optimiser correctly zeroes unhelpful
hardware rates ($\gamma_\text{AD} \to 0$) while preserving teacher
rates, achieving $F_\text{mixed} \approx F_\text{teacher}$ or
slightly better.  This selective behaviour is consistent across all
fidelity-loss levels and represents a robust ``do no harm'' property
of the rate-Krotov framework.


\section{Discussion}

The central finding is that \textbf{all-to-all ZZ dephasing with
optimised rates} provides the best balance of expressivity and
optimisability for the MCWF-Krotov method.  The mechanism is
physically intuitive: ZZ interactions create parity-selective
decoherence that, combined with single-qubit rotations, can implement
\emph{effective} multi-qubit dissipation.  The gate layers provide the
rotational degree of freedom to align the desired and undesired
subspaces with the ZZ eigenstates, while the dissipation suppresses
the orthogonal complement.

The failure of Config~D (ancilla-reset) highlights a fundamental
limitation of the two-stage pipeline: the Krotov update
(Eq.~\ref{eq:krotov_gate}) optimises only gate parameters, while the
Kraus channel parameters remain fixed from Stage~1.  In principle,
the ancilla-reset angles $\boldsymbol\phi$ could be treated as
additional ``gate'' parameters with their own generators.  However,
the non-unitary character of the Kraus map breaks the clean
forward--backward structure of Krotov's method for those specific
parameters.  A hybrid approach combining Krotov updates for gates
with gradient descent for ancilla angles is a promising direction
for future work.

The gap between Config~C ($F = 0.93$) and the ideal teacher tail
($F = 0.997$) suggests room for improvement.  The extended Krotov
method (Section~\ref{sec:rate_results}) addresses one of these avenues
by directly co-optimising dissipation rates during the Krotov loop.
On a mixed (hardware + teacher) basis, the joint optimiser achieves
$F = 0.88$ without any pre-fitting stage, and demonstrates the ability to
selectively retain helpful operators while disabling harmful ones.

\paragraph{Rate-Krotov selective suppression.}
A striking result from the joint optimiser is that \emph{all} hardware
operator rates (AD, dephasing, ZZ) converge to zero across all bases
and all system sizes (4, 6, and 8~qubits; Tables~\ref{tab:rate_krotov}
and~\ref{tab:rate_krotov_scaling}), while teacher-cooling rates are
retained where beneficial.  This confirms that the rate-Krotov gradient
(Eq.~\ref{eq:krotov_rate}) correctly captures the functional utility of
each operator.  The per-layer rate structure---where teacher operators
are active only in later layers---reveals a temporal optimisation
strategy not available to the two-stage pipeline.

\paragraph{Scaling bottleneck: Krotov update magnitude.}
The 6- and 8-qubit experiments (Section~\ref{sec:scaling_rate_krotov})
reveal that the $O(1/\sqrt{d})$ scaling of the Krotov update is the
fundamental barrier.  At $d = 64$ (6q), the joint mode still achieves
$F = 0.55$ on the richest basis, but at $d = 256$ (8q) improvements
are marginal.  This does \emph{not} invalidate the rate-Krotov
extension---rather, it indicates that the rate optimisation is correctly
functioning (selective suppression is consistent across scales) but the
underlying gate-Krotov updates are too small.  Adaptive step-size
schedules, momentum-based updates, or second-order information could
ameliorate this.

The 10-qubit experiment reveals a richer picture.  The MCWF-Krotov
method scales computationally---running in $25$--$77\,\text{s}$ via
tensor contractions and \texttt{Rank1Op} operators---but the
optimisation landscape becomes much harder.  With Krotov updates of
magnitude $O(1/\sqrt{d}) \approx 0.03$ per gate, the pure-gate
optimiser makes essentially zero progress in 200 iterations.
Teacher dissipation breaks this deadlock ($F = 0.515$), confirming
that target-selective dissipation provides the non-unitary ``push''
needed for convergence at scale.  Unoptimised hardware noise, by
contrast, strictly degrades performance.

These findings point to \textbf{rate optimisation} as the critical
bottleneck for scaling.  Stage~1 of the pipeline requires $O(d^2)$
density-matrix propagation per channel per step, making it prohibitive
at $d = 1024$.  Trajectory-based rate gradient estimation, variational
Lindblad learning, or reduced-rank channel parametrisations are
promising paths to extend the full two-stage pipeline beyond 10~qubits.

\paragraph{Q-Alchemy circuit + dissipative tail.}
The experiment of Section~\ref{sec:qalchemy} tests a practically
motivated scenario: a pre-compiled circuit from Q-Alchemy with a
rate-only dissipative tail.  The results reveal both a positive finding
and a fundamental limitation.  Positively, teacher-cooling operators at
fixed initial rates can close up to $2.2\%$ of the fidelity gap (from
$F = 0.882$ to $F = 0.905$) without any gate modification---a ``free''
improvement for any circuit compiler.  The limitation is that the
first-order Krotov rate gradient vanishes structurally for
target-cooling operators (Eq.~\ref{eq:vanishing_gradient}), preventing
the optimiser from tuning the rates.  This suggests that the cooling
effect is purely passive and that more sophisticated parametrisations
(e.g.\ operator angles rather than rates) would be needed for active
optimisation.


\section{Conclusion}

We have demonstrated that the MCWF-Krotov optimal control method can
be transferred to gate-based quantum circuits with interleaved
hardware-constrained dissipation, and further extended to treat
\emph{Lindblad dissipation rates as Krotov control variables}.
A systematic study at 4~qubits and a scalability demonstration at
10~qubits ($d = 1024$) reveals the following:

\begin{enumerate}[nosep]
\item Pure single-qubit channels (amplitude damping, dephasing) provide
  only modest improvements ($F \lesssim 0.67$) because their structure
  is incompatible with the target state.
\item All-to-all ZZ connectivity with pre-fitted rates achieves
  $F = 0.93$ (4~qubits), demonstrating that long-range correlated
  dissipation is a powerful resource when rates are optimised.
\item The \textbf{extended Krotov method for rates}
  (Eq.~\ref{eq:krotov_rate}) correctly identifies operator utility:
  harmful hardware operators are suppressed to $\gamma \to 0$, while
  target-cooling operators are retained with non-zero rates.
\item \textbf{Joint gates+rates Krotov} achieves $F = 0.88$ (4q) on a
  mixed (hardware + teacher) basis without any pre-fitting stage,
  and maintains $F \ge$ baseline across all bases and system sizes
  (4, 6, 8~qubits) by adaptively disabling harmful channels.
\item The \textbf{per-layer rate structure} discovered by joint
  Krotov reveals a temporal strategy: teacher operators are activated
  primarily in later circuit layers, a degree of freedom not available
  to the two-stage pipeline.
\item The \textbf{protective ``selective silence''} mechanism ensures
  that extending the operator basis never hurts joint-mode
  performance: at 8~qubits, 44 fixed-rate operators destroy
  fidelity ($F = 0.08$), but joint Krotov zeros all rates and
  preserves $F = 0.50$.
\item At 10~qubits ($d = 1024$), pure gate optimisation saturates
  at the truncated-circuit baseline ($F = 0.500$), while
  target-selective teacher dissipation achieves $F = 0.515$---the
  only configuration that improves at this scale.
\item Fixed-rate hardware noise strictly degrades fidelity
  (Table~\ref{tab:10q}), reinforcing the necessity of rate
  optimisation---now available within the Krotov framework itself.
\item \textbf{Pre-compiled circuit + dissipative tail}
  (Section~\ref{sec:qalchemy}): appending a teacher-cooling tail to a
  Q-Alchemy circuit closes up to $2.2\%$ of the fidelity gap at
  6~qubits.  The rate-Krotov gradient vanishes structurally for
  target-cooling operators, identifying a fundamental limitation of the
  first-order formulation and motivating operator-level parametrisation.
\end{enumerate}

These results establish that \emph{dissipation expressivity}, quantified
by the connectivity and type of available channels, is a key design
parameter for dissipation-assisted quantum state preparation.  The
extended Krotov formulation provides a principled, gradient-based
mechanism for jointly optimising coherent and dissipative controls within
a single monotonically convergent framework.  The memory-efficient
\texttt{Rank1Op} formulation and batched density-matrix evaluation
enable the MCWF-Krotov framework to handle $d = 1024$ within minutes,
with $O(d)$ memory per trajectory.  Scaling the joint gate+rate Krotov
method to larger systems, potentially combined with reduced-rank
operator bases, is the most pressing direction for future work.


\begin{thebibliography}{9}
\bibitem{goerz2018}
M.~H.~Goerz and K.~Jacobs,
``Efficient optimization of state preparation in quantum networks
using quantum trajectories,''
\emph{arXiv:1801.04382v2} (2018).

\bibitem{dalibard1992}
R.~Dum, P.~Zoller, and H.~Ritsch,
``Monte Carlo simulation of the atomic master equation for
spontaneous emission,''
\emph{Phys.~Rev.~A} \textbf{45}, 4879 (1992).

\bibitem{plenio1998}
M.~B.~Plenio and P.~L.~Knight,
``The quantum-jump approach to dissipative dynamics in quantum optics,''
\emph{Rev.~Mod.~Phys.} \textbf{70}, 101 (1998).

\bibitem{krotov1996}
V.~F.~Krotov,
\emph{Global Methods in Optimal Control} (Dekker, New York, 1996).

\bibitem{palao2003}
J.~P.~Palao and R.~Kosloff,
``Optimal control theory for unitary transformations,''
\emph{Phys.~Rev.~A} \textbf{68}, 062308 (2003).

\bibitem{reich2012}
D.~M.~Reich, M.~Ndong, and C.~P.~Koch,
``Monotonically convergent optimization in quantum control using
Krotov's method,''
\emph{J.~Chem.~Phys.} \textbf{136}, 104103 (2012).

\bibitem{qalchemy}
Q-Alchemy,
``Cloud-based quantum circuit compilation,''
\url{https://www.2q-alchemy.com/} (2024).
\end{thebibliography}

\end{document}
